\section{Related Work}

\paragraph{Automated research systems.}
Recent systems aim to automate parts of the scientific workflow, including ideation, experiment execution, manuscript writing, and self-improvement. Examples include The AI Scientist \citep{lu2024ai}, Agent Laboratory \citep{schmidgall2025agentlaboratory}, and FARS \citep{analemma2025fars}. These systems demonstrate that LLM-based agents can assemble research-like artifacts, but they also raise evaluation questions that are not captured by final-paper appearance alone. \ra\ differs by focusing on off-the-shelf CLI agents under minimal scaffolding and by retaining the full process trace needed for artifact-grounded auditing.

\paragraph{CLI agents and agentic coding systems.}
General-purpose coding agents such as Claude Code, Codex, and Kimi Code operate directly in software repositories and shell environments \citep{anthropic2026opus,anthropic2026claudecode,openai2026codex,moonshot2026kimi,moonshot2026kimicode}. Their ability to edit files, run experiments, and debug failures makes them suitable for end-to-end research pipelines. Coding benchmarks such as SWE-bench \citep{jimenez2024swebench} evaluate issue resolution in real repositories, while \ra\ evaluates a broader research loop that includes problem framing, empirical design, paper writing, and review.

\paragraph{LLM-as-reviewer and paper-review agents.}
LLM-based reviewing systems estimate paper quality, generate review text, or calibrate scores against venue decisions. Stanford Agentic Reviewer provides an ICLR-style score from a submitted manuscript \citep{stanford2026sar}. Related work studies LLM reviewing behavior, reviewer calibration, and biases in automatic evaluation \citep{liang2024llmjudge,zhou2024reviewer2}. \ra\ uses \sar\ as a PDF-only comparison point, but its methodological focus is artifact-grounded review: reviewers inspect the paper together with code, logs, result files, and references.

\paragraph{Reproducibility, artifact evaluation, and scientific integrity.}
Artifact evaluation practices in systems and machine-learning venues emphasize reproducibility, code availability, and the relationship between claims and executable artifacts \citep{pineau2021improving,collberg2016repeatability,acm2020badging}. These practices motivate our evaluation target. A generated paper may be clear and plausible while failing to match its own code or logs. \ra\ adapts the spirit of artifact evaluation to auto-research by making paper-artifact faithfulness a first-class evaluation criterion.
